% ============ Solutions and hints for exercises (Part 3: §7-§9) ============
% 原书 pp.289-291; 原书对 §7 的 Ex.7.1,7.6,7.7、§8 的 Ex.8.5,8.6、§9 的 Ex.9.5,9.7,9.9,9.11 未给解答, 编号按原书跳号.

\section*{\S\ref{sec:7}}

\noindent\textbf{7.2.}\quad For $b \in B$ write $b = y/x$ with $x, y \in A$. Then $y = bx \in Bx \cap A = xA$
(by Theorem~\ref{thm:7.5}, (ii)), so $b \in A$.

\noindent\textbf{7.3.}\quad Write $N \subset M$ for the $A$-submodule generated by $\{m_\lambda\}$; then $B \otimes
(M/N) = 0$, so $M/N = 0$.

\noindent\textbf{7.4.}\quad Set $M = \prod M_\lambda$. It is enough to show that $I \otimes M \to IM$ is injective
for an ideal $I = (a_1, \ldots, a_n)A$ of $A$ (by Theorem~\hyperref[thm:7.6]{6}). Define $f\colon A^n \to A$ by
$f(x_1, \ldots, x_n) = \sum a_ix_i$, and set $K = \Ker f$. Then $0 \to K \to A^n \to A$ is
exact, hence also $0 \to K \otimes M \to M^n \to M$, and if $\sum a_i \otimes \xi_i \in I \otimes M$ satisfies
$\sum a_i\xi_i = 0$ then writing $\xi_{i\lambda}$ for the $\lambda$-th coordinate of $\xi_i \in M$ we have
$\sum a_i\xi_{i\lambda} = 0$ for all $\lambda$, and hence $(\xi_{1\lambda}, \ldots, \xi_{n\lambda}) \in K \otimes M$.
Now since $A$ is Noetherian we can write $K = A\beta_1 + \cdots + A\beta_r$ with $\beta_j
= (b_{1j}, \ldots, b_{nj}) \in K \subset A^n$ for $1 \le j \le r$. Thus we can write $\xi_{i\lambda} = \sum_j b_{ij}\eta_{j\lambda}$
with $\eta_{j\lambda} \in M_\lambda$. Since $\sum_i a_ib_{ij} = 0$, setting $\eta_j = (\eta_{j\lambda})_\lambda \in M$, we get
$\xi_i = \sum_{j=1}^{r} b_{ij}\eta_j$, and $\sum_i a_i \otimes \xi_i = \sum_{i,j} a_ib_{ij} \otimes \eta_j = 0$.

\noindent\textbf{7.5.}\quad Tensor the exact sequence $0 \to A \xrightarrow{a} A$ with $N$ to get the exact sequence
$0 \to N \xrightarrow{a} N$.

\noindent\textbf{7.8.}\quad Tensor product does not commute with infinite direct products. If $\{p_i\}$ is
an infinite set of prime numbers then $\bigcap_i p_i\Z = (0)$, but $\bigcap_i (p_i\Z \otimes \Q) \neq 0$.

\noindent\textbf{7.9.}\quad If $I$ is an ideal of $A$ then $IB \cap A = I$, so that given a chain $I_1 \subset I_2 \subset \cdots$ of
ideals of $A$, $I_1B \subset I_2B \subset \cdots$ eventually terminates, hence so does the given
chain.

\section*{\S\ref{sec:8}}

\noindent\textbf{8.1.}\quad $(I+J)^{2n} \subset I^n + J^n$, so that given $x_1, x_2, \ldots$ such that $x_{n+1} - x_n \in (I+J)^{2n}$
we can write $x_{n+1} - x_n = u_n + v_n$ with $u_n \in I^n$ and $v_n \in J^n$. Thus $\{x_n\}$
converges to $x_1 + \sum u_i + \sum v_i$. Also, $I, J \subset \rad(A)$, so that $I + J \subset \rad(A)$, and
therefore $\bigcap_n (I + J)^n = (0)$.

\noindent\textbf{8.2.}\quad If $\{x_n\}$ satisfies $x_{n+1} - x_n \in J^n$ then there is a limit $x$ for the $I$-adic topology.
For any $i$, taking $m$ large enough we get $x_m - x \in I^i$, so that $x_n - x = x_n - x_m + x_m - x
\in J^n + I^i$, and since by Theorem~\hyperref[thm:8.10]{10}, (i) we have $\bigcap_i (J^n + I^i) = J^n$, we get
$x_n - x \in J^n$, so that $x$ is also a $J$-adic limit.

\noindent\textbf{8.3.}\quad Let $a\hat{A} = x\hat{A}$ and $x = \sum a_i\hat{x}_i$ with $a_i \in a$ and $\hat{x}_i \in \hat{A}$. Let $I$ be an ideal of
definition of $A$. Take $x_i \in A$ such that $x_i - \hat{x}_i \in I\hat{A}$ and set $a = \sum a_ix_i$. Then
$a \in a$ and $a\hat{A} \subset a\hat{A} + I\hat{A}$, so that by NAK, $x\hat{A} = a\hat{A}$, so $a = x\hat{A} \cap A = aA$.

\noindent\textbf{8.4.}\quad $e = (e - X)(e + eX + eX^2 + \cdots)$.

\noindent\textbf{8.7.}\quad $A/m^n$ is Artinian, so that there exists $t(n)$ such that $a_{t(n)} + m^n = a_j + m^n$ for
$j > t(n)$. We can assume that $t(n) < t(n+1) < \cdots$. Supposing that $a_{t(n)} \nsubseteq m^r$ for some
$r$, then we take $a_r \in a_{t(r)} - m^r$, then $a_{r+1} \in a_{t(r+1)}$ such that $a_{r+1} - a_r \in m^r$, and
proceed in the same way taking $a_i \in a_{t(i)}$ such that $a_i - a_{i-1} \in m^{i-1}$ for $i \ge r$. Then
$\lim a_i$ belongs to $\bigcap_\nu a_\nu$, but not to $m^r$, which is a contradiction.

\noindent\textbf{8.8.}\quad This can be done by following the proof of Theorem~\hyperref[thm:8.5]{5}, and replacing the
use of homogeneous polynomials by multihomogeneous polynomials.

\noindent\textbf{8.9.}\quad We can assume that $A$ is a Noetherian local ring with maximal ideal $P$.
There is an $x \neq 0$ such that $xP = 0$, and then since $\bigcap_n P^n = (0)$ there exists $c$
such that $x \notin P^c$. Then if $I \subset P^c$, $I:x = P$.

\noindent\textbf{8.10.}\quad Let $m = (X, Y) \subset k[X, Y]$ and set $A = k[X, Y]_m$; let $q(X) = \sum c_iX^i \in k[[X]]$ be
transcendental over $k(X)$ and set $a_\nu = (X^{\nu+1}, Y - 1 - c_1X - \cdots - c_\nu X^\nu)$.

\section*{\S\ref{sec:9}}

\noindent\textbf{9.1.}\quad $B_p$ is integral over $A_p$, so that any maximal ideal of $B_p$ lies over $pA_p$ and
therefore coincides with $PB_p$. Hence $B_p$ is a local ring, and the elements of $B - P$
are units of $B_p$.

\noindent\textbf{9.2.}\quad $\dim A \le \dim B$ from the going-up theorem, and $\dim B \le \dim A$ from Theorem~\hyperref[thm:9.3]{3}, (ii).

\noindent\textbf{9.3.}\quad Replacing $A$ and $B$ by $A_p$ and $B_p$ we can assume $p$ is maximal; then set
$k = A/p$, so that $B/pB$ is a finite $k$-module, hence an Artinian ring.

\noindent\textbf{9.4.}\quad If $ax^n \in A$ for all $n$ then $A[x]$ is a submodule of the finite $A$-module $a^{-1}A$; if
$A$ is Noetherian then $A[x]$ is also a finite $A$-module.

\noindent\textbf{9.6.}\quad Suppose $f = gh$ with $g, h \in K[X]$ monic. Roots of $g$ are roots of $f$, hence
integral over $A$, and expressing the coefficients of $g$ in terms of the roots, we
have that the coefficients of $g$ are integral over $A$; since $A$ is integrally closed,
$g \in A[X]$, and similarly for $h$.

\noindent\textbf{9.8.}\quad By Theorem~\hyperref[thm:9.3]{3}, (ii).

\noindent\textbf{9.10.}\quad $L[X]$ is a free module over $K[X]$, hence flat; and if $L$ is algebraic over $K$
then $L[X]$ is integral over $K[X]$. The first part follows from the previous
two questions, together with Theorem~\hyperref[thm:9.5]{5}. If $f, g$ have a common factor $\alpha(X)$
in $L[X]$ then set $P = (\alpha(X))$, so that $\htop P = 1$ (it can easily be seen that a
non-zero principal prime ideal in a Noetherian integral domain has height
$1$). Hence $\htop p \le 1$. But $f, g \in p$, so that $\htop p = 1$. There is an irreducible
divisor $h$ of $f$ in $p$, and $p = (h)$, so $h \mid g$.
